% ============================================================================
%  Formulario T-11. Especificaciones técnicas ofertadas (FEP01 p. 62)
%  Subdocumento 4.2 según el Formulario T-21.
%
%  Qué se compra, cuánto y con qué características (Caso, capítulo 11, p. 24,
%  y FEP02, capítulo 8, pp. 18 a 20). Cada componente es un registro:
%
%  \componenteTOnce{
%    componente      = Dispositivo a bordo,
%    producto        = marca y modelo de referencia, o servicio ofertado,
%    caracteristicas = características mínimas (RT-08.01, RT-08.10 a RT-08.12),
%    ubicacion       = lugar de instalación,
%    cantidad        = cantidad total, con repuestos si corresponde,
%    adquiere        = CLIENTE o audIT,
%    ciclo           = ciclo de vida y reposición (RT-08.13),
%    justificacion   = por qué este componente, aquí y en esta cantidad
%  }
%
%  No hay columna de costo. RT-08.10 pide «costo unitario estimado», pero el
%  Artículo 50.2 prohíbe toda cifra que permita inferir la oferta económica:
%  la nota al pie de la tabla remite ese dato al Sobre N.º 3.
% ============================================================================

\DefinirFormulario{T-11}{Especificaciones técnicas ofertadas}{62}
  {Detalle de los componentes de infraestructura, plataforma, licenciamiento
   y hardware especificado, con su ubicación o lugar y su justificación. El
   hardware lo adquiere el CLIENTE y el PROPONENTE debe especificar
   exactamente qué comprar, cuánto y con qué características
   (Caso, capítulo 11, p. 24).}

\newcommand\componenteTOnce[1]{\RegistroNuevo{t11}{#1}}

\ExplSyntaxOn
\NewDocumentEnvironment { formularioTOnce } { O{} }
  { \ReiniciarRegistros { t11 } }
  {
    \IniciarFilas
    \int_step_inline:nn { \CantidadRegistros { t11 } }
      {
        \tl_gput_right:Ne \g_audit_filas_tl
          {
            ##1 &
            \exp_not:N \textbf { \audit_campo:nnn { t11 } {##1} { componente } }
            \exp_not:N \par
            \exp_not:N \audit@rotuloceldas { Adquiere } ~
            \audit_campo:nnn { t11 } {##1} { adquiere } &
            \audit_campo:nnn { t11 } {##1} { producto }
            \exp_not:N \par
            \exp_not:N \audit@rotuloceldas { Características } ~
            \audit_campo:nnn { t11 } {##1} { caracteristicas }
            \exp_not:N \par
            \exp_not:N \audit@rotuloceldas { Ciclo~de~vida } ~
            \audit_campo:nnn { t11 } {##1} { ciclo } &
            \audit_campo:nnn { t11 } {##1} { ubicacion } &
            \audit_campo:nnn { t11 } {##1} { cantidad } &
            \audit_campo:nnn { t11 } {##1} { justificacion } \exp_not:N \\
          }
      }
    \TablaAcumulada [ ancho = completo , fuente = condensada , #1 ]
      { Formulario~T-11.~Especificaciones~técnicas~ofertadas } { t11 }
      { R{6mm} L{28mm} L{50mm} L{22mm} L{17mm} X }
      { N.º & Componente & Producto~o~servicio~ofertado & Ubicación & Cantidad & Justificación \\ \midrule }
    \notaTabla
      {
        Costo~unitario~estimado~(\transv{RT-08.10}{19}):~se~declara~en~el~
        Sobre~N.º~3,~Oferta~Económica,~conforme~a~\art{50.2}{29}.~Esta~oferta~
        técnica~no~contiene~precios,~tarifas~ni~valores~unitarios.
      }
  }
\ExplSyntaxOff

% Rótulo dentro de una celda
\newcommand\audit@rotuloceldas[1]{{\color{audit-secundario}#1.}}

% Esqueleto: filas vacías, con todos los campos por completar
\expandafter\def\csname esqueleto@T-11\endcsname{%
  \begin{formulario}{T-11}
    \begin{formularioTOnce}
      \componenteTOnce{}
      \componenteTOnce{}
      \componenteTOnce{}
    \end{formularioTOnce}
  \end{formulario}}
